% ===== tr-01 块3：第 2 章 pp.11--15（原文页 11--15，正文止于 References）=====
% [原文 p.11]
尽管此情形下判别轨迹具有实余维一，SYZ 镜像构造仍可实施，并如所期望地重新得到环面 Calabi--Yau 三维流形 $K_{\PP^2}$。\footnote{更确切地说，式 (2.6.1) 的 SYZ 镜像是 $K_{\PP^2}$ 中一个光滑超曲面的补集。}
% [Tutor Note: "conic bundle"（二次曲线丛）指式 (2.6.1) 所定义的 $\check X\subset\C^2\times(\C^\times)^2$：对固定的 $(z_1,z_2)$，方程 $uv=\text{常数}$ 是 $\C^2$ 中的一条二次曲线，故 $\check X$ 可看作 $(\C^\times)^2$ 上的二次曲线丛。]

然而，在环面世界之外，尚不清楚如何显式地进行 SYZ 构造。一个主要问题是开 Gromov--Witten 不变量的良定义性。只在少数非环面情形（见 Liu [130] 与 Solomon [161]）我们才有良定义的开 Gromov--Witten 不变量理论。\footnote{不过，Fukaya 最近的工作 [54, 55] 在紧 Calabi--Yau 三维流形上给出了圆盘计数不变量的定义。}

\section{非 Calabi--Yau 情形下的 SYZ}
镜像对称被发现后不久，就通过 Batyrev [10]、Givental [62, 63, 64]、Kontsevich [112]、Hori--Vafa [98] 等许多人的工作被推广到非 Calabi--Yau 情形，尤其是 Fano 流形。与 Calabi--Yau 情形不同，镜像不再由一个流形给出；取而代之的是，它被预言为一对 $(\check X,W)$，其中 $\check X$ 是非紧 Kähler 流形，$W:\check X\to\C$ 是全纯函数。在物理学文献中，这样的一对 $(\check X,W)$ 称为 Landau--Ginzburg 模型，$W$ 称为该模型的超势 [168, 176]。
% [Tutor Note: 非 Calabi--Yau 情形下镜像不再是流形，而是 Landau--Ginzburg 模型 $(\check X,W)$：$\check X$ 非紧 Kähler、$W$ 全纯（超势）；不要把它与紧 Calabi--Yau 情形的"镜像流形 $\check X$"混为一谈。]

自然会问：SYZ 提议能否也推广到这一情形？Auroux [8] 第一个考虑了这个问题，而且他考虑的是一个远为一般的情形，即由紧 Kähler 流形 $X$ 与一个有效反典范除子 $D$ 组成的对 $(X,D)$。$D$ 的定义截面定义了 $X$ 上的一个亚纯体积形式，它只沿 $D$ 有单极点（且在 $X\setminus D$ 上处处非零），从而使我们可以在补集 $X\setminus D$ 上谈论特殊拉格朗日环面纤维化。
% [Tutor Note: 正是"只沿 $D$ 有单极点、在 $X\setminus D$ 上非零"这一点使 $X\setminus D$ 上存在全纯体积形式，从而可以定义特殊拉格朗日子流形；这是 Auroux 推广的关键技术条件。]

假设给定这样的纤维化 $\mu:X\setminus D\to B$。那么我们可尝试运行 SYZ 构造以产生 SYZ 镜像 $\check X$，即考虑对 $(L,\nabla)$ 的模空间（其中 $L$ 是 $\mu$ 的纤维，$\nabla$ 是 $L$ 上的平坦 $U(1)$-联络），然后再作瞬子修正。此外，超势 $W$ 会自然地出现，作为与 Fukaya--Oh--Ohta--Ono 关于 $\mu$ 的拉格朗日纤维的 Floer 复形的障碍链 $m_0$ 相对应的镜像对象。
% [Tutor Note: $m_0$ 是 Fukaya--Oh--Ohta--Ono 理论中的障碍链（obstruction chain），即 Floer 微分之外的第一个 $A_\infty$-结构项；超势 $W$ 正是它的镜像对象，此处不可把 $m_0$ 理解为普通映射或矩阵。]

当 $X$ 是紧环面 Kähler 流形时，$D$ 的一个典范取法是所有环面主除子之并。此时矩映射给出 $X$ 上一个方便的拉格朗日环面纤维化，它有一个很好的性质：在开稠密的环面轨道 $X\setminus D$ 上限制为一个环面丛。在此情形，SYZ 镜像流形 $\check X$ 就由代数环面 $(\C^\times)^n$ 给出，因为我们有一个环面丛，且在镜像流形的构造中没有瞬子修正。所有本质信息都编码在超势 $W$ 中。在 Auroux 的工作之前，是 Cho 与 Oh [39, 41] 首先注意到 $W$ 可以用圆盘计数不变量（或亏格 0 开 Gromov--Witten 不变量）来表达。通过分类 $X$ 中由矩映射纤维界定的所有全纯圆盘，他们得到了 $X$ 为 Fano 情形下 $W$ 的显式公式，这与 Hori--Vafa [98] 用物理论证得到的更早预言一致。这一点后来被 Fukaya--Oh--Ohta--Ono 的工作 [57, 58, 59] 极大地推广到所有紧环面流形。
% [Tutor Note: 环面情形镜像流形恰为 $(\C^\times)^n$ 有两个原因：(1) 在 $X\setminus D$ 上矩映射给出环面丛（半平坦模型）；(2) 镜像流形的构造中没有瞬子修正。全部非平凡信息都进入超势 $W$，这与紧 Calabi--Yau 情形不同。]

% [原文 p.12]
在 [22] 中，环面 Fano 流形的镜像对称被用作一个试验场，以考察 Fourier 型变换（我们称之为 SYZ 变换）在研究镜像对称几何时能有多大用处。对环面 Fano 流形 $X$，我们考虑开稠密的环面轨道 $U_0:=X\setminus D\subset X$，它也是矩映射的拉格朗日环面纤维之并。在辛几何意义下，我们可以写 $U_0=T^*B_0/\Lambda^\vee$，其中 $B$ 是矩多面体，$B_0$ 表示其内部。此时 SYZ 镜像是 $\check X:=TB_0/\Lambda$，它是 $(\C^\times)^n$ 中的一个有界区域。为得到超势 $W$，我们考虑 $X$ 中逐纤维的测地/仿射环路构成的空间
\[
\tilde X:=U_0\times\Lambda\subset\mathcal{L}X,
\]
在 $\tilde X$ 上，我们有一个瞬子修正的辛结构 $\tilde\omega=\omega+\Phi$，其中 $\Phi$ 是亏格 0 开 Gromov--Witten 不变量的生成函数，这些不变量计数由矩映射纤维界定的（虚拟）全纯圆盘。
% [Tutor Note: $\mathcal{L}X$ 表示 $X$ 的自由环路空间（loop space），不是"$\mathcal{L}$ 乘以 $X$"；$\tilde X=U_0\times\Lambda$ 是其中由"逐纤维的测地/仿射环路"构成的子空间，$\Lambda$ 记录环路绕纤维的绕数。]
% [Tutor Note: 注意 $\check X:=TB_0/\Lambda$ 是 $(\C^\times)^n$ 中的有界区域（因为 $B_0$ 是矩多面体的内部），与 p.4 的 $\check X=TB/\Lambda$ 是同一半平坦构造在 $B$ 有边界时的版本。]

随后，通过把半平坦 SYZ 变换 $\mathcal{F}^{\text{semi-flat}}$ 与逐纤维的 Fourier 级数结合起来，构造出一个显式的 SYZ 变换 $\mathcal{F}$，并证明了 $\mathcal{F}$ 把 $X$ 上修正后的辛结构 $\tilde\omega$ 恰好变成镜像 Landau--Ginzburg 模型 $(\check X,W)$ 的全纯体积形式 $e^W\check\Omega$，其中 $W$ 由对 $\Phi$ 作逐纤维 Fourier 变换得到。
% [Tutor Note: 原文此处写"$\tilde\omega$ on $X$"，但 $\tilde\omega=\omega+\Phi$ 是前一段定义在 $\tilde X$ 上的辛结构；正文照原文排作 $X$，读者应按上下文理解为 $\tilde X$。]

此外，$\mathcal{F}$ 诱导了 $X$ 的（小）量子上同调环 $QH^*(X)$ 与 $W$ 的 Jacobian 环 $Jac(W)$ 之间的同构。其证明通过取热带极限，并观察到其上同调对应物对量子乘积有贡献的热带曲线可以由热带圆盘粘接得到；细节见 [22]。这一观察后来被 Gross [74] 推广，用于他通过热带几何对 $\PP^2$ 的大量子上同调所做的镜像对称研究。
% [Tutor Note: $QH^*(X)$ 是（小）量子上同调环，$Jac(W)=\C[x_i]/(\partial W/\partial x_i)$ 是 $W$ 的 Jacobian 环；同构 $QH^*(X)\cong Jac(W)$ 是环同构，而非仅仅向量空间同构。]

至于一般型流形，目前沿 SYZ 视角研究其镜像对称主要有两条途径。一条是 Abouzaid--Auroux--Katzarkov [3] 的工作，他们考虑环面簇 $V$ 中的一个超曲面 $H$，并构造了一个 Landau--Ginzburg 模型，它是 $V\times\C$ 沿 $H\times\{0\}$ 的爆破的 SYZ 镜像。特别地，当 $H$ 是 $V=\PP^1\times\PP^1$ 中一个双次数 $(3,2)$ 多项式的零点集时，他们的构造产生亏格 2 Riemann 曲面的一个镜像，这与 Katzarkov 先前的提议 [108, 106, 157] 一致。

另一条更接近 Gross--Siebert 纲领的途径是 Gross--Katzarkov--Ruddat 的工作 [80]，他们提出一般型簇的镜像是一个带某种消没轮（vanishing cycles）层的可约簇。据推测，这条途径产生的镜像应当给出与 [3] 相同的数据。例如，该可约簇应当是 SYZ 镜像 Landau--Ginzburg 模型的超势的临界轨迹。但这两条途径之间的确切关系仍在研究中。

\section{超越 SYZ}
除了给出镜像对称的优美几何解释之外，SYZ 猜想 [163] 还对许多相关的数学领域产生了深远影响。我们下面简要描述这方面应用的几个例子。

% [原文 p.13]
\para{HMS 经由 SYZ。}正如我们已经看到的，SYZ 猜想基于弦论中 D-膜的思想。回忆：B-膜（即 B-模型中的 D-膜）是复子簇上的凝聚层，而 A-膜（即 A-模型中的 D-膜）是带平坦 $U(1)$-联络的特殊拉格朗日子流形。因此，我们可以把 Kontsevich 的同调镜像对称（HMS）猜想 [111]（它断言 Calabi--Yau 流形 $X$ 的（导出）Fukaya 范畴等价于镜像 $\check X$ 上凝聚层的导出范畴）看作 $X$ 上 A-模型与 $\check X$ 上 B-模型之间同构的一种体现。于是很自然地可以期望：SYZ 提议、特别是 SYZ 变换，可被用来构造实现 HMS 猜想所断言的范畴等价的几何函子。

例如，给定 SYZ 纤维化 $\mu:X\to B$ 的一个拉格朗日截面，它与 $\mu$ 的纤维 $L$ 的交点决定了对偶环面 $L^\vee$ 上的一个平坦 $U(1)$-联络。把这些平坦 $U(1)$-联络粘接起来，就得到对偶纤维化全空间上的一个全纯线丛，而后者在模去量子修正后就是镜像 $\check X$。这个简单的想法最早由 Gross [69, 70] 构想，随后被 Arinkin--Polishchuk [7] 与 Leung--Yau--Zaslow [122] 用来构造 SYZ 变换，并进而用于证明与理解半平坦 Calabi--Yau 情形下的 HMS 猜想。后来，同样的想法也被用来研究环面簇的 HMS 猜想 [1, 2, 46, 48, 47, 16, 24, 40]。
% [Tutor Note: 对偶环面 $L^\vee$ 上的平坦 $U(1)$-联络之所以给出全纯线丛，是因为 $L$ 的每个交点对应一个 $U(1)$-相位；粘接这些相位即得对偶纤维化全空间上的线丛，这是 SYZ 变换在对象层面的几何实现。]

在一些更近的工作 [18, 28, 27, 26, 81] 中，SYZ 变换被用来构造（在对象层面上的）几何 Fourier 型函子，它们对某些环面 Calabi--Yau 流形的例子实现了 HMS 范畴等价，例如 $A_n$-奇点的消解与光滑化的锥形流形；在这些例子中会遇到带奇异纤维、因而具有非平凡量子修正的 SYZ 纤维化。

另一方面，K.-L. Chan、Leung 与 Ma 正在进行的工作 [29, 30] 表明，SYZ 变换也可以用来在态射层面上构造 HMS 等价，至少在半平坦情形如此。最终目标是：对任给的 SYZ 纤维化构造一个典范的几何 Fourier 型函子，实现 HMS 猜想所断言的范畴等价，从而加深我们对 HMS 猜想的几何、乃至整个镜像对称的理解。

\para{Ricci 平坦度量与圆盘计数。}SYZ 论文 [163] 中一个值得注意的观察是：镜像上的 Ricci 平坦度量可以分解为一个半平坦部分（早先已在 Greene--Shapere--Vafa--Yau [68] 中显式写出）与一个瞬子修正部分之和，后者应当来自原 Calabi--Yau 流形中由 SYZ 纤维化的拉格朗日环面纤维所界定的全纯圆盘的贡献。这提示了一种具体而定性描述 Ricci 平坦度量的方式，而这类度量几乎从来不是显式的。

一般而言，这样的定性描述仍然极难获得，因为紧 Calabi--Yau 流形上 SYZ 纤维化的非平凡例子难以找到，且开 Gromov--Witten 理论尚未被很好理解。然而，Gaiotto--Moore--Neitzke 的近期工作 [60, 61] 为超 Kähler 情形带来了新见解。他们提出了一个关于复可积系统全空间上的超 Kähler 度量（最简单的例子是著名的 Ooguri--Vafa 度量 [148]）与 Kontsevich--Soibelman 的墙穿越公式之间关系的新猜想。

% [原文 p.14]
为了稍详细地描述他们的提议，我们考虑一个复可积系统 $\psi:M\to B$，即 $M$ 是全纯辛流形，且 $\psi$ 的纤维是复拉格朗日子流形。更确切地说，Gaiotto、Moore 与 Neitzke 在 [60, 61] 中考察的都是亚纯 Hitchin 系统；在这种情形，完备的超 Kähler 度量最早由 Biquard--Boalch [11] 构造（例如见 [61, 第 4.1 节]）。他们利用了这样一个事实：任何超 Kähler 度量都由其相伴的 扭量空间刻画，因此他们尝试在 $M$ 上构造一族由 $\C^\times$ 参数化的全纯 Darboux 坐标卡，使其满足 Hitchin 等人 [96] 定理的假设。特别地，他们要求这些坐标满足某些墙穿越公式，这些公式描述了坐标在所谓 BPS 射线两侧的不连续性，而 BPS 态的（虚拟）计数在穿过 BPS 射线时发生跳变。
% [Tutor Note: $\psi:M\to B$ 是复可积系统（全纯辛流形 $M$ + 复拉格朗日纤维化），与 SYZ 纤维化 $\mu:X\to B$ 结构类似但处于全纯辛而非 Kähler 范畴；$B$ 在此并非 SYZ 的底流形。]
% [Tutor Note: "BPS rays"（BPS 射线）是 $B$ 中使 BPS 态计数跳变的射线；墙穿越公式正是描述坐标跨过这些射线时的不连续性，这与 Kontsevich--Soibelman 散射图中的墙穿越公式同型。]

这些墙穿越公式原来与 Kontsevich--Soibelman [115] 以及 Gross--Siebert [87] 在构造 Calabi--Yau 流形的环面退化时所用的公式等价（另一方面，它们也与 motivic Donaldson--Thomas 理论 [105, 113] 中的墙穿越公式相同）。考虑到这一点以及 SYZ 猜想，并且注意到 [11] 已通过超 Kähler 旋转构造出亚纯 Hitchin 系统上大量的非紧 SYZ 纤维化例子（其中包括复二维的许多例子，例如引力瞬子、对数 Calabi--Yau 曲面），自然可以期望：Gaiotto、Moore 与 Neitzke 所考虑的那些复可积系统例子上的超 Kähler 度量，可以用全纯圆盘的（虚拟）计数来表达。

这一点对最简单的例子------[17] 中的 Ooguri--Vafa 度量------已经做到。W. Lu 的更近工作 [134, 135] 表明，一般而言，[60, 61] 所研究的亚纯 Hitchin 系统上的 扭量空间与全纯 Darboux 坐标所产生的数据，与运行 Gross--Siebert 纲领 [87] 所需的数据相同，从而说明超 Kähler 度量与热带圆盘计数之间必定存在某种（或许是隐含的）联系。Y.-S. Lin 在其博士论文 [129] 中考虑了椭圆 K3 曲面与热带圆盘计数不变量，并证明他的不变量满足与 [60, 61] 中出现的相同的墙穿越公式。显然，这些 K3 曲面上的超 Kähler 度量也与圆盘计数密切相关。Stoppa 及其合作者的更近工作 [162, 49] 也表明了 motivic Donaldson--Thomas 理论中的墙穿越公式与 Gaiotto--Moore--Neitzke 的构造之间的密切联系。

\para{SYZ 的其它应用。}我们还想提及 SYZ 构造最近的两个出人意料的应用，但不深入细节。

在最近的一个合作项目 [77] 中，Gross--Hacking--Keel 构造了对数 Calabi--Yau 曲面的 SYZ 镜像族；所谓对数 Calabi--Yau 曲面是指一对 $(Y,D)$，其中 $Y$ 是非奇异射影有理曲面，$D\in|-K_Y|$ 是有理曲线组成的循环。他们通过把 [87] 中的构造推广到允许更一般（即更坏）奇点类型的整仿射流形来做到这一点。令人惊奇的是，他们的结果可用来证明 Looijenga [133] 一个 30 年前的猜想，该猜想关乎尖点奇点的可光滑化。

在更近的一篇论文 [78] 中，他们应用自己的构造证明了对数 Calabi--Yau 曲面的 Torelli 定理，该定理最初由 Friedman [52] 于 1984 年猜想。另一方面，他们的构造还与 cluster 簇的理论密切相关，并且他们提出了 Fock--Goncharov 对偶基 [76, 79] 的一个宏大推广。关于这些激动人心的新结果与新进展的精彩阐述，我们建议读者参阅 Gross 与 Siebert 的精彩综述文章 [83]。

% [原文 p.15]
在另一个方向上，SYZ 构造出人意料地被用来构造新的纽结不变量。对 $S^3$ 中的纽结 $K$，其余法丛 $N^*K$ 典范地是余切丛 $T^*S^3$ 中的一个拉格朗日循环。在 [43] 中，Diaconescu--Shende--Vafa 在消解锥形流形 $X:=\mathcal{O}_{\PP^1}(-1)\oplus\mathcal{O}_{\PP^1}(-1)$ 中构造了相应的拉格朗日循环 $L_K$；粗略地说，这是通过把余法丛 $N^*K$ 抬离零截面、并让 $T^*S^3$ 经历锥形跃迁而做到的。他们的构造受到物理学中一个被称为大 $N$ 对偶的神秘现象的启发。
% [Tutor Note: 消解锥形流形 $X:=\mathcal{O}_{\PP^1}(-1)\oplus\mathcal{O}_{\PP^1}(-1)$ 是 $\PP^1$ 上两个次数为 $-1$ 的线丛的直和（全空间），它是非紧 Calabi--Yau 三维流形；注意两个 $\mathcal{O}_{\PP^1}(-1)$ 不可写成 $\mathcal{O}_{\PP^1}(1)$。]
% [Tutor Note: 余法丛 $N^*K\subset T^*S^3$ 是 $K$ 的余法丛，作为余切丛的子流形典范地是拉格朗日循环；把它"抬离零截面"并用锥形跃迁 $T^*S^3\rightsquigarrow X$ 搬到消解锥形流形中，即得 $L_K$。]

在 [6] 中，Aganagic--Vafa 通过对对 $(X,L_K)$ 应用推广的 SYZ 构造定义了一个新的纽结不变量。更确切地说，他们的不变量是对 $(X,L_K)$ 的亏格 0 开 Gromov--Witten 不变量的生成函数。结果表明，所得函数总是一个多项式，并且他们猜想它正是纽结理论中经典 A-多项式的形变 [42]。此外，有人提出了他们的不变量与 $K$ 的接触同调代数 [144] 的增广之间的一个有趣关系。最近的一篇论文 [5] 给出了这一关系的大量证据。

SYZ 猜想的这两个新应用，连同未来还会出现的更多应用，开辟了镜像对称以及数学与物理的许多其他分支的新方向\footnote{一个我们没有提及的有趣故事是 Gukov--Yau--Zaslow [93] 把 SYZ 图景应用于 $G_2$ 流形，其目的是几何地解释 M-理论与杂化弦理论之间的对偶。}；它们都指向未来更优美、更激动人心的研究工作。
% [Tutor Note: 脚注中的 $G_2$ 流形记号是 $G_2$（G 带下标 2，特殊 holonomy 群），不是 $G^2$ 或 $G_2$ 的平方；M-理论、杂化弦理论（heterotic string）保留通用译名。]
